\documentclass[a4paper, serif, 10pt]{moderncv}

%% moderncv
% style and color
\moderncvstyle{banking}
\moderncvcolor{black}

%% page layout/margins
\usepackage[top=2.5cm, bottom=2.5cm, left=1.5cm, right=1.5cm]{geometry}

\nopagenumbers{}

%% Babel config for German language
% ref:
% - https://latex3.github.io/babel/guides/locale-german.html
\usepackage[ngerman]{babel}

%% fontspec
\usepackage{fontspec}

\IfFontExistsTF{Linux Libertine}{
  \setmainfont[Ligatures={TeX, Common}, Numbers=OldStyle]{Linux Libertine}
}{}
\IfFontExistsTF{Linux Libertine O}{
  \setmainfont[Ligatures={TeX, Common}, Numbers=OldStyle]{Linux Libertine O}
}{}

%% CTeX
% CJK support, used to show CJK characters in the resume
%
% - fontset=none: disable builtin fontset but instead set the CJK font manually
% - heading=false: disable ctex heading
% - punct=kaiming: use kaiming punctuations styles for CJK
% - scheme=plain: use plain scheme, do not override \normalsize font size
% - space=auto: space settings for CJK characters
%
% ref:
% - http://ctan.mirrorcatalogs.com/language/chinese/ctex/ctex.pdf
\usepackage[UTF8, heading=false, punct=kaiming, scheme=plain, space=auto]{ctex}

\IfFontExistsTF{Noto Serif CJK SC}{
  \setCJKmainfont{Noto Serif CJK SC}
}{}
\IfFontExistsTF{Noto Sans CJK SC}{
  \setCJKsansfont{Noto Sans CJK SC}
}{}

%% line spacing
\usepackage{setspace}
\setstretch{1.125}

%% URL styling - use normal text instead of monospace
\urlstyle{same}

% Auto-underline all links
% Defer until \begin{document} so hyperref is already loaded
\AtBeginDocument{%
  \let\oldhref\href
  \renewcommand{\href}[2]{\oldhref{#1}{\underline{#2}}}%
}

%% Patch \httplink and \httpslink to handle full URLs
% The original moderncv macros always prepend http:// or https:// to URLs.
% This causes issues when the URL already has a protocol (e.g., https://example.com)
% resulting in malformed URLs like https://https://example.com.
% This patch checks if the URL already contains "://" and uses it directly if so.
% Note: We use \str_set:Nx (x-expansion) to expand macros like \@homepage first.
\ExplSyntaxOn
\str_new:N \g_yamlresume_url_str

\RenewDocumentCommand{\httpslink}{O{}m}{
  \str_set:Nx \g_yamlresume_url_str {#2}
  \str_if_in:NnTF \g_yamlresume_url_str {://}
    { \url{#2} }
    { \url{https://#2} }
}

\RenewDocumentCommand{\httplink}{O{}m}{
  \str_set:Nx \g_yamlresume_url_str {#2}
  \str_if_in:NnTF \g_yamlresume_url_str {://}
    { \url{#2} }
    { \url{http://#2} }
}
\ExplSyntaxOff

\name{Anna Weber}{}
\title{Senior Data Scientistin für Machine Learning \& Predictive Analytics}
\phone[mobile]{+49 170 5557890}
\email{anna.weber@datacraft.de}
\homepage{https://www.datacraft.de/anna-weber}

\address{Musterstraße 12, Berlin, Berlin, Deutschland, 10115}{}{}

\extrainfo{{\small \faLinkedin}\ \href{https://www.linkedin.com/in/anna-weber}{@anna-weber} {} {} {} • {} {} {}
{\small \faGithub}\ \href{https://github.com/annaweber}{@annaweber}}

\begin{document}

\maketitle

\section{Basisinformationen}

\cvline{}{Data Scientistin mit über sechs Jahren Erfahrung in der Entwicklung und Implementierung von Machine-Learning-Modellen und skalierbaren Datenpipelines.
%
\begin{itemize}
\item Expertise in Python, statistischer Analyse und Big-Data-Technologien wie Apache Spark und Kubernetes
\item Nachgewiesene Erfolge bei der Steigerung von Umsatz und Kundenzufriedenheit durch datengetriebene Entscheidungen
\item Erfahrung in der Zusammenarbeit mit cross-funktionalen Teams und der Kommunikation komplexer Ergebnisse an nicht-technische Stakeholder
\item Leidenschaft für die Verbesserung von Modellgenauigkeit und die Automatisierung von ML-Workflows
\end{itemize}}
\section{Ausbildung}

\cventry{Okt. 2015–Sept. 2017}
        {Master, Data Science, Punkte: 1,3}
        {Technische Universität München}
        {\url{https://www.tum.de}}
        {}
        {\begin{itemize}
\item Abschlussarbeit: „Ein Deep-Learning-Ansatz zur Vorhersage des Energiebedarfs in intelligenten Stromnetzen"
\item Vertiefung in Machine Learning, statistischer Datenanalyse und verteilten Systemen
\item Praktische Erfahrung durch Projekte mit TensorFlow, Spark und Cloud-Technologien
\end{itemize}
\textbf{Kurse}: Maschinelles Lernen, Statistische Modellierung, Big Data Infrastrukturen, Datenvisualisierung, Deep Learning, Zeitreihenanalyse}

\cventry{Okt. 2012–Sept. 2015}
        {Bachelor, Informatik, Punkte: 2,0}
        {Universität Heidelberg}
        {\url{https://www.uni-heidelberg.de}}
        {}
        {\begin{itemize}
\item Bachelorarbeit: „Entwicklung einer interaktiven Visualisierung zur explorativen Datenanalyse"
\item Grundlagen in Mathematik, Informatik und Statistik
\end{itemize}
\textbf{Kurse}: Algorithmen und Datenstrukturen, Datenbanken, Softwaretechnik, Statistik}
\section{Berufserfahrung}

\cventry{Jan. 2020–Heute}
        {Senior Data Scientistin}
        {DataCraft GmbH}
        {\url{https://www.datacraft.de}}
        {}
        {\begin{itemize}
\item Leitung der Entwicklung einer ML-Plattform für Kundensegmentierung und personalisierte Empfehlungen
\item Aufbau und Automatisierung von Data-Pipelines auf AWS mit Docker, Kubernetes und Terraform
\item Einführung von MLflow zur Modellverwaltung und Reproduzierbarkeit
\item Enge Zusammenarbeit mit Produktmanagement und Softwareentwicklung zur Integration von Modellen in Produktionssysteme
\end{itemize}
\textbf{Schlüsselwörter}: Machine Learning, Python, Kubernetes, AWS, MLflow}

\cventry{Okt. 2017–Dez. 2019}
        {Data Scientistin}
        {FinanzInsight AG}
        {\url{https://www.finanzinsight.de}}
        {}
        {\begin{itemize}
\item Entwicklung von Risikomodellen zur Betrugserkennung im Zahlungsverkehr
\item Analyse großer Transaktionsdatensätze mit Apache Spark und SQL
\item Erstellung interaktiver Dashboards zur Überwachung von Modell-Performance und Geschäftskennzahlen
\item Präsentation von Ergebnissen vor Fachbereichsleitungen und externen Partnern
\end{itemize}
\textbf{Schlüsselwörter}: Risikomodellierung, Apache Spark, SQL, Tableau}

\cventry{Apr. 2015–Sept. 2016}
        {Werkstudentin Datenanalyse}
        {Analytik Solutions GmbH}
        {\url{https://www.analytiksolutions.de}}
        {}
        {\begin{itemize}
\item Unterstützung des Data-Science-Teams bei der Datenaufbereitung und explorativen Analyse
\item Automatisierung von Reporting-Prozessen mit Python und Excel
\item Durchführung statistischer Auswertungen für Kundenprojekte
\end{itemize}
\textbf{Schlüsselwörter}: Python, Statistik, Excel, Reporting}
\section{Sprachen}

\cvline{Deutsch}{Muttersprache oder zweisprachig \hfill \textbf{Schlüsselwörter}: Muttersprache}
\cvline{Englisch}{Verhandlungssichere Kenntnisse \hfill \textbf{Schlüsselwörter}: TOEFL 108, Verhandlungssicher}
\cvline{Französisch}{Gute Kenntnisse \hfill \textbf{Schlüsselwörter}: Grundkenntnisse}
\section{Fähigkeiten}

\cvline{Machine Learning}{Experte \hfill \textbf{Schlüsselwörter}: Python, TensorFlow, PyTorch, scikit-learn, XGBoost}
\cvline{Datenanalyse}{Experte \hfill \textbf{Schlüsselwörter}: Pandas, NumPy, SQL, Apache Spark, Tableau}
\cvline{Cloud Computing}{Erfahren \hfill \textbf{Schlüsselwörter}: AWS, Azure, Docker, Kubernetes, Terraform}
\cvline{Statistik}{Erfahren \hfill \textbf{Schlüsselwörter}: Hypothesentests, Regressionsanalyse, Clusteranalyse, Zeitreihen}
\section{Auszeichnungen}

\cventry{Sept. 2019}
        {Bitkom}
        {KI-Innovationspreis 2019}
        {}
        {}
        {Auszeichnung für die Arbeit zur Optimierung von Prognosemodellen für den Energiebedarf mittels Deep Learning.}
\section{Zertifikate}

\cventry{Feb. 2021}
        {Amazon Web Services}
        {AWS Certified Machine Learning – Specialty}
        {\url{https://aws.amazon.com/certification/aws-certified-machine-learning-specialty/}}
        {}
        {}
\section{Veröffentlichungen}

\cventry{Juni 2021}
        {Vorhersage von Kundenabwanderung mittels Ensemble-Methoden}
        {Zeitschrift für Praktische Datenwissenschaft}
        {\url{https://www.zfpd.de/publikationen/abwanderung-ensemble}}
        {}
        {Vergleich verschiedener Ensemble-Methoden zur Prognose der Kundenabwanderung im Telekommunikationssektor.}
\section{Projekte}

\cventry{Feb. 2020–Nov. 2020}
        {End-to-End-Machine-Learning-Pipeline zur Identifikation von Kundinnen und Kunden mit hohem Abwanderungsrisiko.}
        {Churn Prediction Platform}
        {\url{https://github.com/annaweber/churn-prediction}}
        {}
        {\begin{itemize}
\item Entwicklung einer skalierbaren ML-Pipeline mit automatisiertem Retraining und Monitoring
\item Integration von Kafka zur Echtzeit-Datenverarbeitung
\item Deployment als Docker-Container auf Kubernetes
\end{itemize}
\textbf{Schlüsselwörter}: Python, MLflow, AWS, Kafka}
\section{Interessen}

\cvline{Datenvisualisierung}{D3.js, ggplot2, Storytelling mit Daten}
\cvline{Bergsteigen}{Alpen, Trekking}
\section{Ehrenamtliche Tätigkeiten}

\cventry{Juni 2019–Heute}
        {Freiwillige Datenexpertin}
        {Code for Germany}
        {\url{https://codefor.de}}
        {}
        {\begin{itemize}
\item Unterstützung gemeinnütziger Organisationen bei der Analyse offener Daten
\item Entwicklung datenbasierter Anwendungen für bürgernahe Projekte
\end{itemize}}

\end{document}